
\head{第十版新核查：第0章20.6--20.7}
\source{EGA IV第一部分，分册印刷第136--153页／上传PDF第135--152页；作者勘误据IV第四部分，印刷第345页／上传PDF第344页}
本轮从（20.6.1）连续译至（20.7.20），共46个原著编号条目，保留50个原编号公式或交换图标签。20.6“非完美性模与特征同态”和20.7“向拓扑环的推广”均已译入；原文随后进入§21。本轮的重点是核与张量的先后次序、复形同调的符号、连续导子的标量环，以及完备化后只能得到的稠密性。

新增正文采用核实后的表述，原文差异在相应脚注注明。下面区分排印或引文修正、补明证明步骤，以及原文已经提出而译文必须保留的限制。\textbf{没有将这些事项计成同样数量的“错误定理”，也没有宣称已经独立证明本轮及历史稿件中的全部结论。}

\section*{本轮处理清单}
\addcontentsline{toc}{section}{本轮处理清单}
\begin{longtable}{@{}p{.18\linewidth}p{.19\linewidth}p{.54\linewidth}@{}}
\toprule 标识 & 原著位置 & 判定与处理\\\midrule\endhead
J10-T01 & 20.6.16 & 次数零直和中的补模下标应为 $K_0$，不是 $K_1$。\\
J10-G02 & 20.6.16--17 & 原映射公式与通常投影识别产生两个负号；补明符号调整，正合列结论不变。\\
J10-T02 & 20.6.25 & 证明末尾的引文改为陈述已给出的18.3.6.4及18.3.6.2。\\
J10-T03 & 20.7.2 & 张量--Hom对应右端的标量环应为 $B$，不是 $C$。\\
J10-C01 & 20.7.8 & 依作者正式勘误，把两处18.4引文改成18.5的相应条目。\\
J10-T04、T05 & 20.7.11、20.7.13 & 同态编号与所用正合列编号纠正；不改命题假设。\\
J10-T06 & 20.7.16 & 无限超越次数及未完备微分模不有限生成的两项断言明写 $n\ge1$。\\
J10-T07 & 20.7.18 & 提升产生的是 $A$-导子，原文误称 $B$-导子；用恒等同态核验。\\
J10-G01、G03 & 20.6.11、20.7.17 & 展开分裂判据、可数逆系统的提升与稠密性论证；不当作对原定理的反例。\\
J10-N01--N04 & 相应条目 & 保留非完美性模的适用边界、普通Hom复形、张量后取核及历史宣告的准确含义。\\
\bottomrule
\end{longtable}

\audit{J10-T07}{提升中的导子到底相对于哪个基环？}{原证明的标量环误写；已改入正文}
\oldauditsrc{（20.7.18），分册印刷第152--153页／上传PDF第151--152页}
原证明令
\[
 L=\Omega^1_{B/A}\widehat\otimes_B C,\qquad E=D_C(L)=C\oplus L,
\]
并以连续 $A$-导子
\[
 D:B\longrightarrow L
\]
给 $E$ 赋予 $B$-代数结构
\[
 f:B\to E,\qquad b\mapsto(v(b),D(b)).
\]
提升给出截面 $g:C\to E$，满足 $gv=f$。原文随即称
$g(c)=(c,D'(c))$ 中的 $D'$ 为“连续 $B$-导子”。这个称呼不能与一般情况下的 $D'v=D$ 并存。

\lead{直接按环同态的定义计算。}
平方零平凡扩张的乘法为
\[
 (c,\ell)(c',\ell')=(cc',c\ell'+c'\ell).
\]
因此 $g$ 的可加性及乘法相容性恰好给出
\[
 D'(c+c')=D'(c)+D'(c'),\qquad
 D'(cc')=cD'(c')+c'D'(c).
\]
又因 $g$ 为 $A$-同态，$D'$ 在 $A$ 的像上为零，故 $D'$ 是 $A$-导子。
条件 $gv=f$ 则给出
\[
 (v(b),D'(v(b)))=(v(b),D(b)),
 \qquad D'v=D.
\]
若把 $D'$ 称为 $B$-导子，则依定义必须有 $D'v=0$，这会不当地强迫 $D=0$。

\lead{具体核对。}
取任意域 $k$，
\[
 A=k,\qquad B=C=k[t],\qquad v=1_{k[t]},
\]
全部赋离散拓扑。这些环可容许且可度量，$v$ 形式平展。此时
\[
 L=k[t]\,dt,\qquad D(t)=dt\ne0.
\]
条件 $D'v=D$ 要求 $D'=D$。而任何 $B$-导子 $C=B\to L$ 恒为零，故不可能是原文所称的 $B$-导子。

正文改成“连续 $A$-导子”，保留 $D'v=D$。随后由完备微分模的表示性质，
$D'=h\widehat d_{C/A}$；在稠密生成子集 $D(B)$ 上，
\[
 h\widehat v_{C/B/A}D=D,
\]
因而 $h\widehat v_{C/B/A}=1_L$，原命题的充分性证明得以继续。这里的修正不否定命题结论。

\audit{J10-G02}{两项复形短正合列的符号核对}{补明识别中的负号；不是定理反例}
\oldauditsrc{（20.6.16）--（20.6.17），分册印刷第142--143页／上传PDF第141--142页}
为避免多重下标遮蔽符号，记
\[
 U=\Omega^1_{A/\Lambda}\otimes_A C,\quad
 V=\Omega^1_{B/\Lambda}\otimes_B C,\quad W=\Omega^1_{C/\Lambda},
 \qquad U\xrightarrow f V\xrightarrow g W.
\]
这里 $f,g$ 是原条目指定的自然映射，不假定 $gf=0$。
三个两项复形分别为
\[
 K_L=(U\xrightarrow fV),\quad K=(U\xrightarrow{gf}W),\quad
 K_R=(V\xrightarrow gW).
\]
给 $K$ 加上可缩复形 $(V\xrightarrow1V)$，得到
\[
 K'_1=U\oplus V,\qquad K'_0=W\oplus V,\qquad
 d'(y,z)=(gf(y),z).
\]
原文所取映射为
\[
 \begin{array}{ll}
 j_1(x)=(x,f(x)),&p_1(y,z)=z-f(y),\\
 j_0(x)=(g(x),x),&p_0(y,z)=g(z)-y.
 \end{array}
\]
先逐一核对：
\[
 d'j_1(x)=(gf(x),f(x))=j_0f(x),
\]
\[
 p_0d'(y,z)=g(z)-gf(y)=gp_1(y,z).
\]
每个 $j_i$ 单射，每个 $p_i$ 满射，且
$\Ker p_1=\operatorname{Graph}(f)$、
$\Ker p_0=\operatorname{Graph}(g)$，所以确实得到逐次分裂的复形短正合列。

\lead{负号从哪里产生。}
用投影到 $K$ 的第一直和因子识别 $H_i(K')\simeq H_i(K)$。
次数一的循环可写为 $(y,0)$，其中 $gf(y)=0$；于是
\[
 p_1(y,0)=-f(y).
\]
因此 $H_1(p)$ 是自然映射 $\Ker(gf)\to\Ker(g)$ 的相反数。
次数零的同调类可用 $(w,0)$ 表示，且
\[
 p_0(w,0)=-w.
\]
因此 $H_0(p)$ 也是自然商映射 $\Coker(gf)\to\Coker(g)$ 的相反数。
原证明省略了这两个符号。它们不破坏正合性，因为对模同态 $a$，
\[
 \Ker(-a)=\Ker a,\qquad\Im(-a)=\Im a.
\]

\lead{边界映射的正号。}
取 $z\in\Ker g$。在 $p_1$ 下取提升 $(0,z)$，则
\[
 d'(0,z)=(0,z)=j_0(z).
\]
故连接同态为
\[
 \Ker g\longrightarrow\Coker f,\qquad z\longmapsto[z],
\]
正是正文（20.6.17.2）的自然映射。将上述两个由 $p$ 诱导的箭头改取相反数，就得到原定理使用的所有自然映射，正合列本身成立。

\lead{伴随的次数下标修正（J10-T01）。}
原文将 $\operatorname{Graph}(g)\subset W\oplus V$ 的补模写成
$K_1(C/A/\Lambda)\oplus\{0\}$。该处在次数零，补模应为
\[
 K_0(C/A/\Lambda)\oplus\{0\}=W\oplus\{0\}.
\]
这是一处对象下标错误，与前面的负号问题分开处理。

\audit{J10-N03}{“先取核再张量”与“先张量再取核”}{原文已有区别；译文不能将两者合并}
\oldauditsrc{（20.6.14.5）、（20.6.21，i），分册印刷第141、144页／上传PDF第140、143页}
原文上标 $C$ 的记号表示
\[
 \Upsilon^C_{B/A/\Lambda}
 =\Ker\bigl(\Omega^1_{A/\Lambda}\otimes_A C
       \to\Omega^1_{B/\Lambda}\otimes_B C\bigr).
\]
它与 $\Upsilon_{B/A/\Lambda}\otimes_B C$ 之间只有典范同态；一般不能直接识别。原文在（20.6.21）已经给出保证识别成立的平坦性条件，本报告不把此事当作原著错误。

\lead{一个非平坦商的计算。}
取特征不为2的域 $k$，
\[
 \Lambda=k,\qquad A=k[x],\qquad B=k[x]/(x^2),\qquad
 C=B/(x)=k .
\]
有
\[
 \Omega^1_{A/k}\otimes_A B=B\,dx,\qquad
 \Omega^1_{B/k}=B\,dx/(2x\,dx)=(B/(x))\,dx .
\]
因此
\[
 \Upsilon_{B/A/k}=(x)\,dx .
\]
由于 $x(x)=0$，该 $B$-模与 $k$ 同构，于是
\[
 \Upsilon_{B/A/k}\otimes_B C\simeq k.
\]
另一方面，张量以后，相关映射成为
\[
 k\,dx\xrightarrow{\ 1\ } k\,dx .
\]
它的核为零，即 $\Upsilon^C_{B/A/k}=0$。故典范映射在此例为
$k\to0$，不是同构。

\lead{原文的两个充分条件。}
若 $C$ 对 $B$ 平坦，张量保持原正合列，结论成立。
若 $\Omega^1_{B/\Lambda}$ 及 $\Omega^1_{B/A}$ 对 $B$ 平坦，
令 $I=\Ker(\Omega^1_{B/\Lambda}\to\Omega^1_{B/A})$，则 $I$ 也平坦。
把原正合列拆为
\[
0\to I\to\Omega^1_{B/\Lambda}\to\Omega^1_{B/A}\to0,
\]
\[
0\to\Upsilon_{B/A/\Lambda}\to\Omega^1_{A/\Lambda}\otimes_A B\to I\to0,
\]
分别与 $C$ 张量；两列的右端平坦，故两次张量均保持左端单射。这才得到所需同构。

\audit{J10-G01}{特征同态表示扩张类时的分裂条件}{补全原文的简略证明，并区分“一个L”与“所有L”}
\oldauditsrc{（20.6.11），分册印刷第140页／上传PDF第139页}
设
\[
\begin{gathered}
 M=\Omega^1_{B/A}\otimes_B C,\qquad N=\Omega^1_{C/A},\qquad d:M\to N,\\
 Y=\Ker d,\qquad I=\Im d,\qquad Q=\Coker d.
\end{gathered}
\]
则原同态是限制映射诱导的
\[
 \alpha_L:\Coker\bigl(\Hom_C(N,L)\to\Hom_C(M,L)\bigr)
 \longrightarrow\Hom_C(Y,L).
\]

若 $L$ 内射，任意 $Y\to L$ 可延拓到 $M$，故 $\alpha_L$ 满射。
若 $a:M\to L$ 在 $Y$ 上为零，它先通过 $M/Y\simeq I$ 分解；
再由内射性把 $I\to L$ 延拓到 $N$，故 $[a]=0$，$\alpha_L$ 单射。

若 $Y$ 是 $M$ 的直和因子且 $N\to Q$ 分裂，则 $I$ 是 $N$ 的直和因子。
同样两次延拓可由相应投影完成，不需要 $L$ 内射。

反过来，假定 $\alpha_L$ 对\emph{每一个} $C$-模 $L$ 都双射。
取 $L=Y$，满射性使 $1_Y$ 延拓为 $M\to Y$，故 $Y\subset M$ 分裂。
再取 $L=I$，商映射 $M\to I$ 在 $Y$ 上为零，单射性给出
$N\to I$，其在 $I$ 上的限制为恒等，故 $I\subset N$ 分裂，
也即 $N\to Q$ 分裂。这给出原文所说的充要条件。

若只对\emph{某一个} $L$ 有同构，不能作这个反推。例如 $L=0$ 时同构自动成立，
并不包含任何分裂信息。正文保留原文“对所有 $L$”的限定。

当 $N,Q$ 均投射时，$N\to Q$ 分裂，所以 $I$ 为投射模；
$M\to I$ 也分裂，于是 $Y$ 为 $M$ 的直和因子。这说明原定理最后一句的条件为什么足够。

\audit{J10-T03}{连续张量--Hom对应的标量环}{原式右端的C应为B}
\oldauditsrc{（20.7.2）证明，分册印刷第148页／上传PDF第147页}
该处使用的准确对应是
\[
 \HomCont_C(\Omega^1_{B/A}\otimes_B C,L)
 \simeq \HomCont_B(\Omega^1_{B/A},L).
\]
从左到右限制为 $m\mapsto h(m\otimes1)$；
从右到左把 $a:\Omega^1_{B/A}\to L$ 送到
\[
 m\otimes c\longmapsto c\,a(m).
\]
这两项操作互逆，连续性由张量积拓扑的泛性质给出。

原文右端写成 $\HomCont_C(\Omega^1_{B/A},L)$。
一般的 $\Omega^1_{B/A}$ 只有给定的 $B$-模结构，没有给定的 $C$-模结构，
所以该式一般尚未定义。例如 $A=k,B=k[t],C=k[t,s]$ 时，环同态的方向
$B\to C$ 不能给 $\Omega^1_{B/k}=k[t]\,dt$ 自动赋予 $C$-模结构。
正文修正下标，而不是补造一个额外的标量作用。

\audit{J10-C01}{作者对20.7.8的正式引文勘误}{依原作者勘误落实，不作为新发现}
\oldauditsrc{勘误与增补清单三，A类，分册印刷第345页／上传PDF第344页}
作者明确指定：
\begin{quote}
在（$0_{\mathrm{IV}}$，20.7.8）中，第149页倒数第9行，将（18.4.3）改为（18.5.4）；
倒数第4行，将（18.4.1）改为（18.5.1）。
\end{quote}
这里讨论的是拓扑扩张类，因此应引用18.5的定义及结论。
正文采用修正后的引文，脚注注明原写法；书末有关勘误的选译也收入这一项。
这不表示清单三全文已经译入。

\audit{J10-T02、J10-T04、J10-T05}{三处内部引文的核对}{对象与编号对照；不更改数学假设}
（20.6.25）证明末尾用（18.3.4.2）指代 $u^1,v^1$，
而命题陈述本身已分别给出（18.3.6.4）、（18.3.6.2）。
正文采用这两个明确出处，并用（20.6.6.4）说明与复形构造相符。

（20.7.11）证明所讨论的是
\[
 \Exalcotop_{B/A}(C,L)\to\HomCont_C(\Upsilon_{C/B/A},L),
\]
它是（20.7.10.1），不是原文所引的（20.7.8.1）。
后者的两端都是扩张类模，与这里的目标不同。

（20.7.13）将拓扑情形化为直极限的原依据应为（20.6.25）的导子及扩张类正合列，
不是（20.6.23）的非完美性模正合列。正文作相应修正。
三项均为核对实际源对象、目标对象和证明用途后完成的引用修正。

\audit{J10-T06}{形式幂级数例子的零变量边界}{两项附带断言须排除n=0}
\oldauditsrc{（20.7.16），分册印刷第152页／上传PDF第151页}
原文考虑有限域扩张 $k'/k$ 和
$B=k'[[T_1,\ldots,T_n]]$。其完备微分模有限生成的说明仍适用；
但后面关于分式域具有无限超越次数、在特征零时未完备微分模不有限生成的断言，
须把 $n\ge1$ 写明。

当 $n=0$ 时，$B=k'$，其分式域仍为 $k'$，在 $k$ 上超越次数为零。
若 $\operatorname{char}k=0$，有限扩张 $k'/k$ 可分。对任意 $a\in k'$，
最小多项式 $f$ 满足
\[
 0=d(f(a))=f'(a)\,da,\qquad f'(a)\ne0,
\]
从而 $da=0$，故 $\Omega^1_{k'/k}=0$，当然是有限型模。
这仅否定把原文后两项断言延伸到零变量的读法，不否定 $n\ge1$ 的例子。

正文在后两项断言前加入 $n\ge1$，脚注说明边界；
并未把这个条件加到前面的（20.7.15）或其初始适用性上。
本轮没有把“非零变量时分式域的无限超越次数”改写成一个新的独立证明。

\audit{J10-G03}{完备化后的正合性不能被稠密性替换}{展开原文允许的结论，保留其限制}
\oldauditsrc{（20.7.17）、（20.7.20），分册印刷第152--153页／上传PDF第151--152页}
原文明确说完备化后的序列不必正合；可度量情形给出的结论是末箭头满射，
以及前一箭头的像在中间核中\emph{稠密}。译文保持这一区别。

取可数的相容开理想商，令
\[
 F_k=\Omega^1_{B_k/A}\otimes_{B_k}C_k,\quad
 G_k=\Omega^1_{C_k/A},\quad H_k=\Omega^1_{C_k/B_k}.
\]
每一级有 $F_k\to G_k\to H_k\to0$；转移映射是满射。
令 $I_k=\Im(F_k\to G_k)$，则 $I_{k+1}\to I_k$ 也满射：
先在 $F_k$ 中表示一个像元素，再沿 $F_{k+1}\to F_k$ 提升即可。

\lead{末箭头的满射性。}
给定相容族 $(h_k)\in\varprojlim H_k$。
先取 $g_0\in G_0$ 映到 $h_0$。
若已选 $g_k$，任取 $h_{k+1}$ 的提升 $\widetilde g_{k+1}$；
其在 $G_k$ 的像与 $g_k$ 之差属于 $I_k$。
利用 $I_{k+1}\to I_k$ 满射，把这个差提升并从
$\widetilde g_{k+1}$ 中减去，得到同时映到 $g_k$ 和 $h_{k+1}$ 的 $g_{k+1}$。
归纳得到相容族，证明 $\varprojlim G_k\to\varprojlim H_k$ 满射。

\lead{中间像的稠密性。}
设 $g\in\Ker(\varprojlim G_k\to\varprojlim H_k)$，固定 $k$。
则 $g_k\in I_k$，可取 $f_k\in F_k$ 映到它。
利用满转移映射逐级提升并向低级投影，得到相容族 $f\in\varprojlim F_k$。
$f$ 的像与 $g$ 在第 $k$ 个离散商中相同，故该像在中间核中稠密。

这个论证并未使对每个 $k$ 选取的 $f$ 都成为同一个精确原像。
一般还缺少控制 $\Ker(F_k\to I_k)$ 的条件，不能将上式直接加强为正合。
“稠密不等于相等”甚至在最熟悉的拓扑群中也可见：
整数在 $p$-进整数环中稠密，但二者不是同一个集合。
这里只用此例说明逻辑区别，不把它冒充（20.7.17）的环同态反例。

\audit{J10-N01、J10-N02、J10-N04}{本轮保留的三个适用边界}{解释性注释，不增加原文断言}
\lead{非完美性模为零并不等价于一般环的形式光滑。}
取 $A=B=k$、$C=k[\epsilon]/(\epsilon^2)$。
因为 $\Omega^1_{B/A}=0$，所以 $\Upsilon_{C/B/A}=0$。
然而令 $E=k[t]/(t^3)$、$J=(t^2)$，则 $J^2=0$；
同态 $C\to E/J$，$\epsilon\mapsto\bar t$ 无法提升到 $E$：
任何提升必须把 $\epsilon$ 送到 $t+a t^2$，其平方在 $E$ 中为非零的 $t^2$。
故 $C/k$ 不形式光滑。原文（20.6.2）仅说形式光滑蕴涵非完美性模为零；
（20.6.3）中域扩张的等价命题没有因此被否定。

\lead{普通Hom复形不被静默替换为导出Hom。}
（20.6.6）中的
$H^i(K_\bullet,L)=H^i(\Hom_C(K_\bullet,L))$
按原文定义是普通Hom复形的上同调。
本轮没有把它写成任意情形的 $\operatorname{Ext}$ 群，
也没有默认 $K_\bullet$ 的各项投射。

\lead{历史宣告与已经给出的证明分开。}
（20.6.26）宣布两项复形与余法丛及后续对偶性理论的关系，
并在其所列假设下说明相应同调的识别。
本轮保留其历史内容与证明详略，没有仅凭“同调相同”补称“一般同伦等价”，
也没有用现代余切复形理论替代原文。
（20.7.11）的一般拓扑论证仍沿原文的直极限方式叙述；
本报告没有把有限离散层面的核对宣称为该条所有拓扑细节的独立审定。

\section*{本轮范围与可追踪性}
\addcontentsline{toc}{section}{本轮范围与可追踪性}
本版新增的数学正文为20.6和20.7，共四份分文件源码。
既有数学正文文件不因本轮而整体重写；上一版已标出的H8-P01等问题仍按原状态保留。
本轮所作引文修正、下标修正及补证在正文脚注和本报告中对应列出。
原著目录及已有索引中的页码仍指原出版物，不冒充新合订本页码。

下一连续接续位置为§21“特征 $p$ 的环中的微分”的章首说明及（21.1.1）。
本轮只为核定接续点而读取21.1的开头，没有将它计入已译正文。
第0章§§21--23、第IV章主体及第V卷预备笔记余文尚未连续译完。

技术编译、原条目和公式编号覆盖核对、渲染检查的实际结果保存在工程目录
\texttt{checks/}；这些检查不代替全部证明的形式化验证或出版级逐句审校。
